%%%%%%%%%%%%%%%%%%%%%%%%%%%%%%%%%%%%%%%%%%%%%%%%%%%%%%%%%%%%%%%%%%%%%%%%%%%%%
%%%
%%% File: utthesis2.doc, version 2.0jab, February 2002
%%%
%%% Based on: utthesis.doc, version 2.0, January 1995
%%% =============================================
%%% Copyright (c) 1995 by Dinesh Das.  All rights reserved.
%%% This file is free and can be modified or distributed as long as
%%% you meet the following conditions:
%%%
%%% (1) This copyright notice is kept intact on all modified copies.
%%% (2) If you modify this file, you MUST NOT use the original file name.
%%%
%%% This file contains a template that can be used with the package
%%% utthesis.sty and LaTeX2e to produce a thesis that meets the requirements
%%% of the Graduate School of The University of Texas at Austin.
%%%
%%% All of the commands defined by utthesis.sty have default values (see
%%% the file utthesis.sty for these values).  Thus, theoretically, you
%%% don't need to define values for any of them; you can run this file
%%% through LaTeX2e and produce an acceptable thesis, without any text.
%%% However, you probably want to set at least some of the macros (like
%%% \thesisauthor).  In that case, replace "..." with appropriate values,
%%% and uncomment the line (by removing the leading %'s).
%%%
%%%%%%%%%%%%%%%%%%%%%%%%%%%%%%%%%%%%%%%%%%%%%%%%%%%%%%%%%%%%%%%%%%%%%%%%%%%%%

%%%%%%%%%%%%%%%%%%%%%%%%%%%%%%%%%%%%%%%%%%%%%%%%%%%%%%%%%%%%%%%%%%%%%%%%%%%%%
%%%
%%
%% This file, and the corresponding tcdthesis.sty the accompanied it, have
%% been modified for the M.Sc. styles used in Trinity College, Dublin
%%
%%
%%%%%%%%%%%%%%%%%%%%%%%%%%%%%%%%%%%%%%%%%%%%%%%%%%%%%%%%%%%%%%%%%%%%%%%%%%%%%
\documentclass[a4paper, 12pt, oneside]{report}         %% LaTeX2e document.
\usepackage {tcdthesis}              %% Preamble.

\mastersthesis                       %% Uncomment one of these; if you don't
% \phdthesis                         %% use either, the default is \phdthesis.

\thesisdraft                         %% Uncomment this if you want a draft
                                     %% version; this will print a timestamp
                                     %% on each page of your thesis.

\leftchapter                         %% Uncomment one of these if you want
% \centerchapter                     %% left-justified, centered or
% \rightchapter                      %% right-justified chapter headings.
                                     %% Chapter headings includes the
                                     %% Contents, Acknowledgments, Lists
                                     %% of Tables and Figures and the Vita.
                                     %% The default is \centerchapter.

% \singlespace                       %% Uncomment one of these if you want
\oneandhalfspace                     %% single-spacing, space-and-a-half
% \doublespace                       %% or double-spacing; the default is
                                     %% \oneandhalfspace, which is the
                                     %% minimum spacing accepted by the
                                     %% Graduate School.

\renewcommand{\thesisauthor}{Charles O'Malley}                %% Your official TCD name.
\renewcommand{\thesismonth}{April}                   %% Your month of graduation.
\renewcommand{\thesisyear}{2026}                      %% Your year of graduation.
\renewcommand{\thesistitle}{Enhancing Model Checker Test Generation with Functional Programming}          %% The title of your thesis; use mixed-case.
\renewcommand{\thesisauthorpreviousdegrees}{, BAI}    %% Your previous degrees, abbreviated; separate multiple degrees by commas.
\renewcommand{\thesissupervisor}{Andrew Butterfield}            %% Your thesis supervisor; use mixed-case and don't use any titles or degrees.
% \renewcommand{\thesiscosupervisor}{}                %% Your PhD. thesis co-supervisor; if any.

% \renewcommand{\thesiscommitteemembera}{}
% \renewcommand{\thesiscommitteememberb}{}
% \renewcommand{\thesiscommitteememberc}{}
% \renewcommand{\thesiscommitteememberd}{}
% \renewcommand{\thesiscommitteemembere}{}
% \renewcommand{\thesiscommitteememberf}{}
% \renewcommand{\thesiscommitteememberg}{}
% \renewcommand{\thesiscommitteememberh}{}
% \renewcommand{\thesiscommitteememberi}{}


\renewcommand{\thesisauthoraddress}{...}

\renewcommand{\thesisdedication}{...}     %% Your dedication, if you have one; use "\\" for linebreaks.


%%%%%%%%%%%%%%%%%%%%%%%%%%%%%%%%%%%%%%%%%%%%%%%%%%%%%%%%%%%%%%%%%%%%%%%%%%%%%
%%%
%%% The following commands are all optional, but useful if your requirements
%%% are different from the default values in tcdthesis.sty.  To use them,
%%% simply uncomment (remove the leading %) the line(s).

\renewcommand{\thesisdegree}{Master of Science in Computer Science}
                                     %% default is "DOCTOR OF PHILOSOPHY"
                                     %% for \phdthesis or "MASTER OF ARTS"
                                     %% for \mastersthesis.  Provide the
                                     %% correct FULL OFFICIAL name of
                                     %% the degree.
\renewcommand{\thesisdegreestream}{ (Computer Engineering)}
                                     %% Default is empty. This is used on
                                     %% the title page of the thesis.

\renewcommand{\thesisdegreeabbreviation}{MAI}
                                     %% Use this if you also use the above
                                     %% command; provide the OFFICIAL
                                     %% abbreviation of your thesis degree.
\renewcommand{\thesistype}{Thesis}    %% Use this ONLY if your thesis type
                                     %% is NOT "Thesis" for \phdthesis
                                     %% or \mastersthesis.
                                     %% Provide the OFFICIAL type of the
                                     %% thesis; use mixed-case.

%%%
%%%%%%%%%%%%%%%%%%%%%%%%%%%%%%%%%%%%%%%%%%%%%%%%%%%%%%%%%%%%%%%%%%%%%%%%%%%%%

\usepackage{graphicx,color}
\usepackage{anysize}
\usepackage{amsmath}
\usepackage{natbib}
\usepackage{caption}
\usepackage{hyperref}
\usepackage{listings}
\usepackage{float}
\hypersetup{hidelinks}

% ============================================================================
% REVISION PLAN FOR NEXT DRAFT
% Overall target:
% - Add roughly 12-15 pages of technical substance, aiming for about 40-45 total
%   pages including front matter and references. That is enough to address the
%   current thin spots without forcing the dissertation all the way to 60 pages.
% - Put almost all new length into Design, Implementation, and Evaluation.
%   Avoid padding the Introduction or chapter map.
%
% Chapter targets:
% - Introduction: +0.5 to 0.75 page. Make the bridge strategy explicit as your
%   own contribution and briefly note that the Haskell work was self-taught.
% - State of the Art: +2 to 2.5 pages. Add more model-checking-to-test-generation
%   references, especially OS/embedded/system-level work, and connect the
%   multilingual boundary discussion more directly to later design choices.
% - Design: +2 to 2.5 pages. Add an end-to-end running example and explain why
%   the subprocess bridge was chosen over deeper integration options.
% - Implementation: +5 to 6 pages. This is the main place to grow: show Python-
%   side changes, bridge message flow, one full translation example, more Haskell
%   code, and shorter better-contained listings.
% - Evaluation: +2 to 2.5 pages. Spell out the experiment procedure, saved trace
%   replay, diff classification, and limits of text-based comparison.
% - Conclusions/Future Work: +1 to 1.5 pages. Rewrite the parity sentence cleanly
%   and explain what pushing the bridge outward by porting spin2test would entail.
%
% Supervisor PDF annotations mapped to actions:
% - p8 "Nice summary" -> keep the Introduction opening concise; do not add bulk here.
% - p9 "bridge approach was your idea" -> add 0.25-0.5 page making that design
%   contribution explicit in Introduction and Design.
% - p9 "Looks good..." -> Dissertation Structure section is fine; no extra length needed.
% - p11 "Any other interesting papers..." -> add 2-3 more closely related references in
%   State of the Art.
% - p11 "Nice" -> Functional-programming framing is working; keep it concise.
% - p12 "You should elaborate..." -> cash out the interface-boundary claims later with
%   concrete bridge/design/implementation detail.
% - p12 "Nice" -> integration-options paragraph is worth keeping as context.
% - p13 "point out that you taught yourself Haskell..." -> add a brief note on this in
%   Introduction or Implementation to frame scope and learning curve honestly.
% - p16 "Good." -> keep the protocol sketch, but surround it with a fuller walkthrough.
% - p17 "focus on the nitty-gritty of HOW" -> trim repeated motivation in Implementation
%   and replace it with mechanics and examples.
% - p18 "say a lot more about HOW the bridge is implemented" -> add message format,
%   Python changes, process lifecycle, and an example exchange.
% - p19 "ensure listings don't get broken..." -> shorten listings, split long ones,
%   and keep only the lines needed for the argument.
% - p21 "Fix these!" -> clean listing flow and wrapping around the state-transition examples.
% - p21 "You need more details here..." -> add a running example from annotation line
%   to generated C, with more Haskell code and at least one translation walkthrough.
% - p24 "Nice example" -> keep the brace-escaping bug as the anchor case study and
%   expand the explanation around it slightly.
% - p25 highlight with no text -> likely rewrite the behavioural-preservation sentence
%   in Conclusions so the claim is stated cleanly.
% - p26 "replace the Python code... by Haskell" -> expand Future Work with the concrete
%   work needed to port spin2test and move the language boundary outward.
% - p27 "references to MC-based testgen for operating systems" -> add these references
%   in State of the Art rather than only in the bibliography.
% ============================================================================

%%------------------------------------------------
%% Listing macros
%%------------------------------------------------
%% Examples for the commands in the document below
%%
%% includecode:
%% \includecode{caption for table of listings}{caption for reader}{filename}
%% - includes a file with code and adds a caption that should describe the code in some detail and a shorter caption for the table of listings
\newcommand{\includecode}[4]{\lstinputlisting[floatplacement=H, caption={[#1]#2}, captionpos=b, frame=single, label={#3}]{#4}}
\lstset{
  basicstyle=\ttfamily\small,
  breaklines=true,
  columns=fullflexible,
  keepspaces=true,
  showstringspaces=false,
  captionpos=b,
  frame=single
}

%%------------------------------------------------
%% Image macros
%%------------------------------------------------

%% includescalefigure:
%% \includescalefigure{label}{short caption}{long caption}{scale}{filename}
%% - includes a figure with a given label, a short caption for the table of contents and a longer caption that describes the figure in some detail and a scale factor 'scale'
\newcommand{\includescalefigure}[5]{
\begin{figure}[htb]
\centering
\includegraphics[width=#4\linewidth]{#5}
\captionsetup{width=.8\linewidth} 
\caption[#2]{#3}
\label{#1}
\end{figure}
}

%% includefigure:
%% \includefigure{label}{short caption}{long caption}{filename}
%% - includes a figure with a given label, a short caption for the table of contents and a longer caption that describes the figure in some detail
\newcommand{\includefigure}[4]{
\begin{figure}[htb]
\centering
\includegraphics{#4}
\captionsetup{width=.8\linewidth} 
\caption[#2]{#3}
\label{#1}
\end{figure}
}





\begin{document}                                  %% BEGIN THE DOCUMENT

\thesistitlepage                                  %% Generate the title page.

%\hypersetup{pageanchor=false}
%\thesisdeclarationpage                            %% Generate the declaration page.

%\thesispermissionpage                             %% Generate the copyright permission page
%\hypersetup{pageanchor=true}

\begin{thesisabstract}                            %% the abstract for your thesis
This dissertation explores whether the refinement stage of an existing SPIN/Promela-based RTEMS test-generation pipeline can be moved from Python to Haskell without changing the behaviour of the generated tests. The work focuses on refine command, a stateful part of the system that turns annotated SPIN traces into C test code. Instead of replacing the whole pipeline at once, the project uses a bridge-based approach so that the Haskell version can be introduced gradually and checked against the original Python implementation. The port eventually moved away from line-by-line translation and towards a more Haskell-native design based on explicit state and smaller transformations. The two versions were then compared across the seven active models in the repository. This process found one real bug in the Haskell output, traced to template brace handling, which was fixed. After that, the remaining differences were limited to harmless formatting. Overall, the project shows that the refiner can be ported in a behaviour-preserving way.
\end{thesisabstract}

%\thesisdedicationpage                            %% Generate the dedication page.

\begin{thesisacknowledgments}                     %% Use this to write your
  Thank you Mum \& Dad.                           %% acknowledgments; it can be anything
\end{thesisacknowledgments}                       %% allowed in LaTeX2e par-mode.
  
  
\tableofcontents                                  %% Generate table of contents.
\listoftables                                     %% Uncomment this to generate list of tables.
\listoffigures                                    %% Uncomment this to generate list of figures.

%%
%% Include thesis chapters here...
%%
\chapter{Introduction}
\label{chap:Introduction}

% REVISION NOTE (target +0.5 to 0.75 page): do not pad the opening summary.
% Add only two short pieces of genuinely useful material in this chapter:
% 1) make it explicit that the bridge-based migration strategy was your own design
%    contribution;
% 2) briefly note that the Haskell side of the work was self-taught, which helps
%    frame the scope and effort honestly.

% SUPERVISOR NOTE (p8): "Nice summary" -> keep this opening tight rather than
% stretching it for length.
This chapter introduces the dissertation, motivates the decision to port part of the RTEMS test-generation pipeline from Python to Haskell, and explains the structure of the document. It first describes the technical context of the project, then states the aim and scope of the work, and finally shows how the remaining chapters support the overall argument.

\section{Motivation}

The repository used in this project contains a SPIN/Promela-based workflow for generating RTEMS validation tests. Promela models are explored with SPIN, the resulting annotated traces are refined into concrete test actions, and those actions are emitted as C test files. The stage studied here is the refinement step, implemented historically in Python. That step is small enough to replace in isolation, but central enough that any change to it can affect the entire pipeline.

The motivation for the project came from the structure of the existing refiner rather than from a failure of the generated tests. The Python version already generated useful RTEMS tests, but it accumulated mutable state, branching logic, and template expansion in one place. That made it harder to reason about the code, and it made larger structural changes risky. Haskell was therefore investigated not as a way to generate different tests, but as a way to re-express the same refinement logic with clearer state handling and stronger separation between stages.

\section{Aim and Scope}

The main aim of the dissertation is to determine whether the \texttt{refine\_command} stage can be ported from Python to Haskell while preserving the behaviour of the existing generator. In practical terms, that means the Haskell version should accept the same annotated SPIN output and produce equivalent test code for the existing RTEMS models.

The scope of the work is deliberately narrow. The dissertation does not attempt to redesign the full \texttt{testbuilder} workflow, replace RTEMS itself, or prove that Haskell is universally better than Python for test generation. Instead, it focuses on one difficult, stateful part of the toolchain and studies how a gradual port can be carried out without losing confidence in the output.

The main contributions of the dissertation are:
% SUPERVISOR NOTE (p9): make it explicit here that the bridge approach was your
% idea, not just an implementation convenience inherited from elsewhere.
% LENGTH TARGET (+0.25 to 0.5 page): add 1 short paragraph before or after this
% list explaining why the bridge strategy is a contribution in its own right.
\begin{itemize}
  \item a bridge-based strategy for introducing a Haskell refiner without discarding the existing Python pipeline;
  \item an explicit-state Haskell implementation of the refinement stage;
  \item an empirical comparison of Python and Haskell generated outputs across the seven active Promela models in the repository.
\end{itemize}

\section{Dissertation Structure}

% SUPERVISOR NOTE (p9): "Looks good..." -> leave this section concise; do not add
% length here unless a later chapter is restructured enough to require a new map.

Chapter~\ref{chap:StateOfTheArt} reviews the background needed for the project, including model-checker-based test generation, functional programming, and multilingual system design. Chapter~\ref{chap:Design} defines the problem more precisely and explains the design decisions behind the bridge-based port. Chapter~\ref{chap:Implementation} describes how the refiner was moved from Python to Haskell and where the implementation changed from line-by-line translation to a more Haskell-native design. Chapter~\ref{chap:Evaluation} evaluates the result by comparing generated outputs and discussing the remaining limitations of the workflow. Chapter~\ref{chap:Conclusions} closes the dissertation and identifies the most useful directions for further work.

\chapter{State of the Art}

At the beginning of each chapter, a description should introduce the reader to the content of the chapter. The description should explain to the reader the layout of the chapter, the contribution that the chapter makes to the overall dissertation and the contribution of the individual sections towards the overall chapter.

From the perspective of this document forming part of your degree, this chapter should demonstrate to the reader your knowledge of the area of your dissertation project. It should present your knowledge in a coherent and detailed form. The reader should understand that you have in-depth knowledge of the area of the dissertation without being overloaded with information.

\section{Background}

A section on the background of the dissertation should provide the reader with an introduction to existing technologies and concepts that form the basis of the
work presented in the dissertation.

\subsection{Functional Programming}
  \href{https://www.cs.kent.ac.uk/people/staff/dat/miranda/whyfp90.pdf}{Why Functional Programming Matters} [1]
    This paper details the benefits of functional programming. Specifically it is said that they have greater capacity for modular design and that the paradigm of lazy evaluation in functional languages allow for simpler code. The reasons detailed within drove the motivation to convert from Python to Haskell  

    The reason given for why functional languages have greater capacity for modular design is through their ability to take functions as arguments to other functions.
    Side effects not being allowed in functional languages create a system where every function has a single well defined output. functions are inputted to other functions rather than have a function do multiple jobs. this allows for the easy modification of lower level functions and higher order functions do not need to worry how the lower levels work, they can be treated as black boxes. when all functions have very specific jobs, they can be reused far more often than a function performing multiple jobs, hence improving modularity.

    Lazy evaluation is a powerful principle. values are only computed when needed. infinite series can be generated and the selectors will compute only the values needed. additionally this separation of generation and selection further improves the modularity of the system. this is a very useful property of these languages for this project in particular, as for this test generation system it allows for generation of a very large number of tests without of running out of memory.
    An example of lazy evaluation being helpful is in evaluating the newton-raphson approximation of square roots. the process is an inherently infinite convergent algorithm where each approximation is better than the last. an imperative langauge would cut off the loop when two approximation differ by no more than some predefined value. Lazy evaluation however allows for the infininite series to be mapped, and then taken to any level of precision once a function calls for it.

    Hughes' arguments were influential in justifying the decision to change part of the codebase to be functional. the increased modularity would aid extending the codebase
    
\vspace{25}
\href{https://www.microsoft.com/en-us/research/wp-content/uploads/2016/07/history.pdf}{A History of Haskell} 
clearly gives a way to write about haskell and its history background and relevance/use to a project like this, more justificaiton for its use

Use it to explain why you chose Haskell:

     Strong static typing and type classes ensure your language backends are complete at compile time.
     Purity and modularity make it easier to reason about transformations from SPIN output to abstract test actions.

Haskell was created at a time where there was a dozen non-strict, functional programming languages in circulation. A conference was was held and it was decided that the lack of a common language was hampering widespread adoption of these class of languages. to remedy this, the committe decided on the creation of a language to unite the space.
the specific principles decided upon at this committee are responsible for its utility to the project at hand. The benefits talked about above still apply due to its nature as a funcitonal programming language but additionally, haskell in particular has its own advantages such as

Purity by default
    other funcitonal languages (Lisps, Scheme, Clojure), have purity as style, side effects are allowed to be used freely. 
    in haskell side effects have to programmed explicitly as functions are pure by default.

    this gives the advantages that functions are easier to reason about. just from the type signature you can  know that the function can throw exception or access the outside world by standard input output.

    it makes it easier to refactor, as when changing code you dont have to worry about affecting mutable varialbles or changing state in a subtle way.

Monads
    monads are the mechanism haskell uses to control side effects and they 


Type classes and ad hoc polymorphism
    type classes are a powerful mechanism in haskell in that they allow the programmer to define an abstract interface with mutliple concrete implementations. For this project this manifests as an abstract "test action" with concrete implementations in various different languages.

    Type classes are flexible in what fields are allowable. in addition haskell enforces completeness in type classes at compile time. this means that for all fields defined for a type class they must be implemetnted



    


\subsection{Multi-Language Systems}
    https://chapering.github.io/pubs/tse24haoran.pdf

    analyses stack overflow posts to answer what developers find hard about multingual development. I used it to learn about the methods of multingual development and solutions to challenges. it details that language interfacing is the biggest challenge faced. there is two method of interfacing, explicit and implicit. explicit interfacing involves the language invoking functions from the second language in the style of same language imports. implicit interfacing involves indirect methods of communicating eg message queues which is what I am using due to the fact that the explicit interface implementations for Haskell to python were limited as will be talked about later.

the paper also suggest more specific solutions for problems. One thing that the developers found tripped developers up a lot was data mismatches in function calls. the solution was to be explicit with all types to ensure no discrepancies    
     https://dl.acm.org/doi/abs/10.5555/3370272.3370280?download=true



     this one doesn't actually let you read it so find another one 
     two papers that talk about the benefits and challenges of multi language systems

\subsection{automated test generation using model checkers}
https://openresearch.surrey.ac.uk/esploro/outputs/conferencePresentation/A-model-checking-based-test-case/99516972702346/filesAndLinks?index=0
this is looking like a paper that does the same sorta thing that andrew butterfields thing does but like way earlier

tjis paper looks at 
Business Process Execution Language, a language that is used to describe web service composition behaviour.
due to its use of concurrency and hierarchy, it was difficult to verify manually. so the author came to a similar solution to ourselves and applied model checkers to automatically create test suites. their approach contrasts ours in theat they use NuSMV instead of Promela



How to tie it in: 

     Use it as a more general example of “model checking as test generator”, not just for your RTEMS context.
     You can say: just as they use a model checker to derive meaningful test sequences for web services, you use SPIN to generate test sequences for RTEMS, then map those to concrete test code in C/Ada/Rust.
     


          
\section{Closely-Related Work}

Work in research areas tends to address a number of specific aspects. Ideally, the discussion of published research should focus on the aspects that have been addressed by various publications - and not a discussion of the individual publications.

For example, if the topic would be a discussion of work on programming languages, the subsections of the related work could be discussions of object orientation and its realisation in various languages or the use of lambda functions by these languages.

\item[\href{https://publications.scss.tcd.ie/theses/diss/2022/TCD-SCSS-DISSERTATION-2022-016.pdf}{SPIN/Promela-Based Test Generation Framework for RTEMS Barrier Manager}] [2]
    This dissertation details previous work on the system I am currently working on. The author's contributions included improving the existing test generation framework, successfully modelling the barrier manager in RTEMS with Promela, and generating a 100\% coverage test suite for the barrier manager. This paper provided significant insight into how the system operates.

    The author identified RTEMS as a good target for formal modelling and test generation due to the fact that RTEMS as a codebase is large, having been worked on for over 30 years. as well as this, due to the fact that RTEMS was originally designed for uniprocessor systems, the test coverage for RTEMS on multi processor systems was less complete. Seeing this they saw it as a  target for the application of Promela/Spin modelling. They modelled the software in Promela and then had Spin output situations where the system fails. this is where the author steps in and created a tool that takes the output, refines it using a custom dictionary that matches specific output phrases and matches them to test code, creating a complete test file for the manager.

    Reading this I came to understand the significance of the RTEMS operating system and the importance of ensuring adequate testing. reading through it gave me the background needed to understand where the test suite generation excelled at and where it could be extended. the system was created to create tests in C. Looking at the codebase it can be seen that in many places, c specific test code is hard coded in. In view of this, my project will look to allow for greater flexibility of test code output by allowing for a language other than c to be used.

https://github.com/ddfisher/HaPy
an iniatllay very promising library that seeks to perform the exact function i needed to make this integration work, calling haskell function from python as a normal importable module. This seemed ideal until further reading of the github reveals that only basic data types are supported, so dicts and the like would not be possible to be exported with this 
// tho note to myself to look at later, do i even need to export the ref dicts, as far as i can remember they are just used as inputs, must ask the ai


https://github.com/nh2/call-haskell-from-anything
this uses an interesting system of wrapping up functions arguments and return values into MessagePacks to export via the C FFI built into haskell, seeks to improve on the base usage of the FFI by simplifying it and extending it to work for any language with an implementation of MessagePacks
\subsection{Aspect \#1}


\subsection{Aspect \#2}


 






    Connect “Why Functional Programming Matters” (Hughes) to your test generation: 

     You already have this, but you can strengthen the link:
         Hughes’ argument about modularity and higher-order functions aligns nicely with your design where:
             Lower-level functions parse YAML, handle refinement, and produce abstract test actions.
             Higher-level functions orchestrate the overall pipeline, treating those stages as black boxes.
             
         Lazy evaluation can be mentioned in the context of generating large test suites without exhausting memory (as you already do in your draft).
         




\section{Summary}

Summarize the chapter and present a comparison of the projects that you reviewed.

\begin{table}[!h]
\begin{center}
	\begin{tabular}{|l|c|c|} 
	\hline
 	\bf  & \bf Aspect \#1  & \bf Aspect \#2 \\
  	\hline
	Row 1 & Item 1 & Item 2 \\
	Row 2 & Item 1 & Item 2 \\
	Row 3 & Item 1 & Item 2 \\
	Row 4 & Item 1 & Item 2 \\
	\hline
	\end{tabular}
\end{center}
\caption[Comparison of Closely-Related Projects]{Caption that explains the table to the reader}	
\label{tab:SummaryProjects}
\end{table}

\chapter{Design}
\label{chap:Design}

% REVISION NOTE (target +2 to 2.5 pages): this chapter should cash out the bridge
% idea with one concrete end-to-end example and a clearer argument for why the
% language boundary sits where it does.

This chapter turns the background of the previous chapter into concrete design decisions. It defines the specific problem within the existing RTEMS toolchain, identifies the constraints that any replacement must satisfy, and explains the bridge-based design used to introduce Haskell gradually rather than by one large swap.

\section{Problem Formulation}
\label{sec:ProblemFormulation}

The problem addressed by this dissertation is not that the existing generator failed to produce useful RTEMS tests. The problem is that one central stage of the pipeline, the Python refiner, became difficult to restructure safely as the project grew. The \texttt{testbuilder} workflow already had working Promela models, a working SPIN-based exploration stage, and working C test generation. Changing the refiner therefore meant changing the part of the system where the risk of regression was highest.

The existing Python refiner combines several kinds of work in one place. It parses annotation lines, tracks context such as whether refinement is currently inside a \texttt{STRUCT} or \texttt{SEQ}, performs lookups in a refinement dictionary, and accumulates generated code fragments. This arrangement works, but it makes it harder to reason locally about how one line of input affects the next state of the translation.

Three properties of the original implementation were especially important. First, context is stored as mutable object state. Second, the main refinement function handles many annotation forms in one place. Third, template expansion and code accumulation are interleaved with control flow. Each of these is manageable in isolation, but together they make structural changes riskier than they need to be.

\subsection{Identified Challenges}

The first challenge was behavioural preservation. The repository already contained working models and a working Python refiner. A replacement could only be justified if it produced the same observable output for the existing model set.

The second challenge was integration. The refiner is not a standalone tool. It is called from the existing Python pipeline and depends on YAML refinement dictionaries, language templates, and per-model header fragments. Any design therefore had to fit around those existing inputs rather than invent a new workflow.

The third challenge was scope control. Replacing the entire pipeline at once would have made it difficult to distinguish language-porting issues from unrelated tooling issues. A narrower design was needed so that the hardest stateful component could be rewritten first and checked continuously against the original behaviour.

\subsection{Proposed Work}

The design adopted in this dissertation is a staged port of \texttt{refine\_command}. The Haskell implementation keeps the same external role as the Python one: it receives tokenised annotation lines, updates its refinement state, and eventually renders a test body. The surrounding Python code is retained so that existing configuration files, model assets, and generation scripts can continue to drive the process.

The central design decision is therefore not only to use Haskell, but to introduce it behind a bridge that preserves the historical Python-facing interface. That keeps the migration incremental. The new implementation can be checked with the same inputs as the old one, and output differences can be detected before they spread into the rest of the workflow.

the bridge was placed at the boundary between \texttt{spin2test.py} and \texttt{refine\_command.py} for good reason. when \texttt{spin2test.py} hands work to \texttt{refine\_command.py}, most of the surrounding setup has already been done, the correct RTEMS model files have been found, the \texttt{.spn} ouput has been split into individual spin trace lines nad the YAML refinement dictionaries specific to the RTEMS model have already been loaded. This leaves \texttt{refine\_command.py} with a more singular and specialised job. it only has to recieve configuration settings at the beginning, process the spin trace lines one by one, and then return the completed test scripts. putthing the bridge here isolates the stateful, pattern matching part of the pipeline while keeping the function of the ported file more consolidated by not also having it handle other parts of the workflow.
this design decision complements the decision to use the bridge to gradually port \texttt{refine\_command.py}. this is because when comparing the output of the new file with the original, if theres a difference then i can be more sure that it has to do with the test conversion pipeline and not with some auxiliary helper function.

Two more aggressive and involved alternatives were considered for this bridge program but were ultimately not chosen.
Firstly a full rewrite at the outset would have moved the bridge location outwards, meaning that more of the system could be interfaced with using haskell, but it would also have linked  \texttt{refine\_command.py} changes with \texttt{spin2test.py}, file orchestration and template loading. as mentioned before this would have made regressions trickier to pinpoint, having a wider area to check. 

Secondly, a Foreign Function Interface (FFI) is another appraoch to the mediating between the haskell and non haskell portions of the code, this time having haskell be called as a compiled library from haskell, like the equivalent of a native library.
this way keeps the refinment inside of one process, as opposed to have the bridge process and the refinement sub process. It wasnt chosen however as it creates a more brittle boundary where things like function signatures, data layout and memory ownership have to be decided exactly. additonally it makes debugging less transparent as the haskell function is called directly by the python code. the stdin/out message pipeline contrasts this by being quite flexible  with its command response system and being extremely visible for debugging, as every message is published to console. these reasons explain why this appraoch was chosen

% SUPERVISOR NOTE (p9): this is a second good place to underline that the bridge
% strategy was your design choice. Add ~0.5 page comparing it with the two obvious
% alternatives: full rewrite first, or tighter FFI-style integration.

\section{Overview of the Design}
\label{sec:OverviewOfDesign}

% LENGTH TARGET (+1 to 1.25 pages): add a running example that starts with one
% concrete annotation line and follows it through the bridge to the emitted C fragment.
% That gives the reader a mental model before Chapter 4 gets into implementation detail.

The overall design separates the system into three logical parts: SPIN produces annotated traces from Promela models, the bridge passes those annotations and configuration data to the Haskell refiner, and the existing Python-side tooling assembles the final test files around the rendered body. This decomposition isolates the port to the refinement stage while leaving the rest of the pipeline stable.

In the repository as it stands now the direction of control remains Python centric.
\texttt{spin2test.py} reads the generated Spin trace file, keeps only lines
that start with \texttt{@@@}, splits those lines into pieces for the test code templates, loads the
refinement dictionary from YAML, and constructs the 
\texttt{command} object. That object name is preserved so the programs that call for it do not have to change, but the implementation behind it is now the
bridge in \texttt{combined\_haskell\_bridge.py}. The bridge starts the
Haskell executable in interactive mode and copies the methods that
the existing Python pipeline expects, including
\texttt{setupLanguage}, \texttt{refineSPINLine\_main}, and
\texttt{test\_body}.

A line-oriented protocol was used between Python and Haskell because it made the interface explicit and easy to inspect. The bridge sends an initial JSON configuration, then a sequence of refinement commands, and finally requests the rendered test body. The protocol is simple enough to log and replay, which proved useful during debugging.

\begin{lstlisting}[caption={Simplified bridge exchange between Python and Haskell}, label={lst:bridge-protocol}]
INIT JSON:{...}
RESPONSE:SUCCESS:INIT JSON received
COMMAND:setupLanguage
RESPONSE:SUCCESS:setupLanguage
COMMAND:refineSPINLine:["4", "CALL", "WaitForSuspend", "2", "resumeRC"]
RESPONSE:SUCCESS:refineSPINLine
COMMAND:testBody
RESPONSE:SUCCESS:testBody:BEGIN
...
RESPONSE:SUCCESS:testBody:END
\end{lstlisting}

Here is a concrete example to help  show why this appraoch works. In the
task-manager model, \texttt{spin2test.py} can encounter the spin 
line \texttt{["4", "CALL", "WaitForSuspend", "2", "resumeRC"]}
while scanning the generated trace. Python does not interpret that
line beyond tokenising it, i.e. reducing every part of the trace line into a separate piece. The token array is sent over the
bridge, and the Haskell refiner then performs the stateful work. It begins by
recording process ID 4, then applying any switch or annotation handling
required by the runtime flags, in this case its "CALL" which doesnt require any switching, then looking up \texttt{WaitForSuspend} in
the refinement dictionary, and formatting the associated template with
the arguments \texttt{2} and \texttt{resumeRC}. 

\begin{lstlisting}[language={}, caption={Refinement template for \texttt{WaitForSuspend}}, label={lst:waitforsuspend-template}]
WaitForSuspend: |
  T_log( T_NORMAL, "Waiting for Task(%d) to Suspend", {0}, {1} );
  do {{
    {1} = ( *ctx->t_isSuspend )( {0} ? taskID[{0}] : 0xffffffff );
    rtems_task_wake_after( 100 );
  }} while ({1} != RTEMS_ALREADY_SUSPENDED);
\end{lstlisting}
here is the template that is called up by searching for "WaitForSuspend" in the refinement dictionaries. the spots where theres {number} are where the tokens are slotted in

\begin{lstlisting}[language=C, caption={Generated C after template substitution}, label={lst:waitforsuspend-generated}]
T_log( T_NORMAL, "Waiting for Task(%d) to Suspend", 2, resumeRC );
do {
  resumeRC = ( *ctx->t_isSuspend )( 2 ? taskID[2] : 0xffffffff );
  rtems_task_wake_after( 100 );
} while (resumeRC != RTEMS_ALREADY_SUSPENDED);
\end{lstlisting}

and here you can see what it looks like when the values from the spin line are subsituted in.

The important design point is that Python does not need to know how that fragment is built;
it only preserves the boundary and passes the same observable inputs
that the old refiner used.

The state after that command remains inside the Haskell process. The
generated lines are accumulated in memory as part of the runtime state
rather than being written to disk immediately. Only when Python later
calls \texttt{test\_body()} does the bridge send
\texttt{COMMAND:testBody}, receive the text body between
\texttt{BEGIN} and \texttt{END} markers, and then assemble the final
output file. It does this by writing the predefined preamble and postamble text, the rendered body, and the run fragment. This separation is why the bridge can
stay narrow: Python keeps responsibility for file-system
orchestration, while Haskell owns the stateful translation from
token array stream to generated C.

JSON plus line-oriented commands were sufficient because only two
payload shapes had to cross the boundary: one configuration object at
startup and one token array per annotation line. Everything else is a
fixed command or a return of the body text. A richer RPC framework would have
added complexity without improving the one property the dissertation
needed most, namely a boundary that could be logged, replayed, and
compared directly against the historical Python behaviour.
% SUPERVISOR NOTE (p16): "Good." Keep this protocol sketch, but add a fuller
% walkthrough around it: who sends each message, what state exists before/after it,
% and why JSON plus line commands were sufficient for this project.

This design also made differential checking a key focus of the project rather than an afterthought. Because the Haskell refiner could be exercised with the same traces as the Python original, text diffs on generated tests became the main feedback loop. That in turn made it possible to move from a line-by-line translation strategy to a more Haskell-native design without losing confidence in the observable output.

Against the two main alternatives, the subprocess bridge performed
best on the criteria that mattered most for this dissertation. It preserved the
historical Python API because callers still interacted with a familiar
\texttt{command}-style object. It was easier to debug because every
request and response could be inspected as plain text. And it supported
debugging checking because the same annotation traces could be replayed
through both implementations with minimal surrounding change. Its main
cost was explicit message passing, but that cost was acceptable because
the boundary was intentionally narrow.

% LENGTH TARGET (+0.75 to 1 page): compare the adopted subprocess bridge with the
% deeper alternatives mentioned in Chapter 2 using explicit criteria: build friction,
% debuggability, preservation of the Python API, and ease of regression checking.

\section{Summary}
\label{sec:SummaryDesign}

In summary, the design of the dissertation is conservative at the system boundary and experimental inside the refiner itself. The boundary stays stable so that behaviour can be checked continuously, while the implementation behind that boundary is free to adopt a clearer Haskell structure.
\chapter{Implementation}
\label{chap:Implementation}
\lstset{language=Python}
\captionsetup{width=.8\linewidth}

% REVISION NOTE (target +5 to 6 pages): most of the honest extra length belongs in
% this chapter. Reduce repeated WHY and replace it with HOW: Python-side changes,
% bridge lifecycle, state transitions, concrete translation examples, and shorter,
% better-focused listings.

This chapter describes how the design of the previous chapter was realised in the repository. It explains the staged transition from the original Python refiner to the Haskell version, the supporting bridge and diff-based workflow used to keep the port testable, and the point at which the implementation stopped being a direct translation and became a more Haskell-native rewrite.

\section{Overview of the Solution}

% SUPERVISOR NOTE (p17): this opening is starting to repeat the WHAT and WHY.
% Trim repetition here and spend the recovered space lower down on mechanics,
% examples, and concrete code-path explanation.

The implementation followed a gradual transition from the original Python refiner to a Haskell refiner while keeping the rest of the pipeline recognisable. The main files involved were \texttt{refine\_command.py}, \texttt{spin2test.py}, \texttt{combined\_haskell\_bridge.py}, and the new \texttt{refine\_command.hs}. The objective was not to replace everything at once, but to replace the refinement stage in a way that could be checked end to end against the old behaviour.

A full rewrite from the start would have been risky. The refiner is small enough to look manageable, but it contains many state-dependent branches and many points where template expansion affects the final test file. Reaching the end of a full rewrite before comparing outputs would therefore have made it difficult to localise errors. The staged approach reduced that risk by keeping differential comparison in the loop throughout the port.

The testing and debugging effort was directed toward functional equivalence with the Python generator. The Python output already matched the expected RTEMS test structure, so the main implementation question was whether the Haskell version could preserve that observable behaviour while using a clearer internal structure.

\section{The Implementation}

\subsection{The Bridge}

% SUPERVISOR NOTE (p18): add 1.5 to 2 pages here on HOW the bridge works in practice.
% Show the Python entry point, the message format, process startup/shutdown,
% error handling, and one short logged exchange from INIT through testBody.
% Also explain what changed on the Python side to preserve the historical command
% interface while delegating the core refinement work to Haskell.

During the transition, a bridge layer preserved the historical Python-facing \texttt{command} interface while moving the refinement work behind a subprocess boundary. Python remained responsible for orchestrating the wider pipeline, but Haskell could now be exercised inside that pipeline without forcing an immediate rewrite of every caller.

Communication was carried entirely through standard input and output. The bridge sends configuration data to the Haskell process, forwards annotation lines one at a time, and finally requests the rendered test body. This arrangement is deliberately plain, but it has two practical benefits: the interface is explicit, and the same messages can be logged or replayed when something goes wrong.

\subsection{The Diff Script}

% SUPERVISOR NOTE (p19): keep listings short and stop them breaking across pages.
% If necessary, replace full-file listings with excerpted fragments that only show
% the mismatch being discussed, and move secondary detail into prose.

A constant part of the porting process was checking the generated test files against the ones produced by the original system. At first this was done manually with an online text comparison tool, but that approach became unworkable once a single change could affect many generated files. A small local diff script therefore became part of the working process.

\includecode{Sample of command line output}{Example output from the diff workflow showing a generated file diverging from the reference output. In practice, even a single extra line was worth investigating because small formatting mismatches could indicate a deeper error in template expansion.}{lst:snippet}{commandline.py}

The diff-based approach was intentionally strict. It reported trivial whitespace differences as well as genuinely important mismatches. That occasionally produced noise, but it was still useful because it made it harder for a small divergence to go unnoticed and later turn into a larger behavioural discrepancy. In practice, most debugging started from a diff, then moved backwards through the relevant refinement dictionary entry and the code path that emitted it.

Unit tests were also added around the Haskell refiner, but they were less decisive than the end-to-end diff checks. Unit tests were useful for protocol handling, configuration loading, and small local behaviours. They were less effective for the interactions that actually mattered most, where several pieces of state and several refinement templates combined to produce one incorrect line of generated code.

\subsection{Refine Command}

% SUPERVISOR NOTE (p21): add 2.5 to 3 pages here with a running example. Start from
% one concrete annotation line or short sequence, show the relevant Python template or
% dictionary entry, show the Haskell state transition(s), and end with the emitted C.
% This is the cleanest way to add useful length without padding.

This was the main file to be converted. Its job is to take annotated SPIN output and refine it into concrete test code by matching annotation lines against entries in a model-specific refinement dictionary. Those entries are small templates of C code with placeholders for names, values, and identifiers taken from the trace.

At the beginning of the port, the Python code was translated into Haskell as closely as possible. The intention was to keep behaviour stable by keeping the structure stable. What became clear quite quickly, however, was that this strategy made the Haskell version look like Python written in another syntax. It preserved the awkward parts of the original without gaining much clarity from the new language.

The turning point came in state handling. The program has to remember whether the current annotation is inside a structure, inside a sequence, or part of ordinary top-level output. The Python version manages this with mutable object fields that are updated as each line is processed. That behaviour can be copied in Haskell, but copying it directly would have reproduced the same difficulty of following how state changes over time.

The more Haskell-native approach was to make the state explicit. Instead of mutating fields on one long-lived object, each stage receives a runtime state, transforms it, and returns the next one. At that point the port stopped being a line-by-line translation and became a Haskell implementation that happened to preserve the same outputs. The aim from then on was not structural similarity to the Python source, but behavioural equivalence of the generated tests.

This change is visible in the following snippets.

\paragraph{Python state stored and mutated on the object}
\begin{lstlisting}[language=Python, caption={Python state stored and mutated on the object}, label={lst:python-state-mutation}]
self.inSTRUCT = False
self.inSEQ = False
self.seqForSEQ = ""
self.inName = ""

case [pid, "SEQ", name]:
    self.inSEQ = True
    self.seqForSEQ = ""
    self.inName = name

case [pid, "SCALAR", "_", value]:
    self.seqForSEQ = self.seqForSEQ + " " + value
\end{lstlisting}

In the Python version the state lives on the command object. As each SPIN line is read these values are updated in place, and later branches depend on whatever earlier branches left behind.

\paragraph{Haskell state made explicit}
\begin{lstlisting}[language=Haskell, caption={Explicit runtime state in the Haskell refiner}, label={lst:haskell-runtime-state}]
data RuntimeState = RuntimeState
  { rtInSTRUCT :: Bool
  , rtInSEQ :: Bool
  , rtSeqForSEQ :: Text
  , rtInName :: Text
  , rtDefCode :: [Text]
  , rtDeclCode :: [Text]
  , rtTestCodes :: [(Int, [Text])]
  , rtCurrentId :: Int
  , rtSwitchNo :: Int
  }
\end{lstlisting}

The same information is still present, but it is now grouped into one explicit record. Rather than hidden mutation on an object, each function takes a state and returns the next one.

\paragraph{One SPIN line producing the next state}
% SUPERVISOR NOTE (p21): "Fix these!" Review the line wrapping and caption flow around
% Listings 4.4 and 4.5. If they still crowd the page, split them or shorten them.
\begin{lstlisting}[language=Haskell, caption={Processing one SPIN line by returning a new state}, label={lst:haskell-state-pipeline}]
refineSPINLineMain :: [Text] -> RuntimeState -> RuntimeState
refineSPINLineMain ln rt0 =
  let rt1 = collectPIdsTokens ln rt0
      rt2 = if outputSWITCH (rtFlags rt1) then switchIfRequired ln rt1 else rt1
      rt3 = if annoteComments (rtFlags rt2) then logSPINLine ln rt2 else rt2
  in refineSPINLine ln rt3
\end{lstlisting}

This is the clearest expression of the change in approach. Instead of one long-lived object being updated in place, each stage produces a fresh version of the runtime state before the next stage runs.

\paragraph{Sequence handling as an explicit transition}
\begin{lstlisting}[language=Haskell, caption={Sequence state handled as an explicit transition}, label={lst:haskell-sequence-state}]
_pidTxt:"SEQ":name:[] ->
  rt { rtInSEQ = True, rtSeqForSEQ = "", rtInName = name }

_pidTxt:"SCALAR":"_":value:[] ->
  rt { rtSeqForSEQ = rtSeqForSEQ rt <> " " <> value }
\end{lstlisting}

Sequence handling was where this change felt most useful. In this context a sequence is an ordered run of values collected between \texttt{SEQ} and \texttt{END} and later emitted through a matching \texttt{\_SEQ} template. In the Haskell version, entering that mode and accumulating those values are both visible directly in the returned state.

% LENGTH TARGET (+1 to 1.5 pages): after the state snippets, add one end-to-end
% translation walkthrough with intermediate runtime-state snapshots. That will satisfy
% the request for more Haskell detail better than adding more isolated code fragments.

\section{Summary}

The implementation was successful because it kept verification in the loop while the internal structure changed. The bridge kept the surrounding Python pipeline stable, the diff script exposed regressions quickly, and the Haskell refiner could then be reorganised into an explicit-state design without losing confidence in the generated output.


\chapter{Evaluation}
\label{chap:Evaluation}

this chapter evaluates the work against the original aim of the dissertation. because the original python generator was already capable of producing satisfactory RTEMS tests, the question here is not whether the haskell port creates a better test suite, but whether it preserves the existing one while fitting into the same workflow.

the evaluation therefore focuses on two things. the first is output parity between the python and haskell refine command across the seven models listed in the repository: barrier-mgr, chains, event-mgr, msg-mgr, sem-mgr, task-mgr, and proto-sem. the second is workflow compatibility, meaning that the haskell version can be used in the existing pipeline without manual fixes. a second experiment on ease of extension had originally been planned, but was not completed in time and is therefore treated as discussion rather than evidence.

\section{Experiments}

for each of the seven models, test code is generated once with the original python refine command and once with the haskell version. the generated files are then compared using the diff script. any differences are examined and classified as either semantic differences, which change the meaning of the generated test, or non-semantic differences, such as formatting.

the run is counted as successful only if no manual intervention is required and any remaining differences do not affect the meaning of the generated test. this makes the evaluation narrow, but it matches the central risk of the project, which is regression in generated output rather than raw performance.


\section{Results}

this gives a clear success criterion for the dissertation. if the haskell refiner can be substituted into the existing workflow and still generate equivalent tests for the full model set, then the main aim of the project has been met. this does not by itself prove that the new code is easier to extend. that claim is supported more weakly by the implementation itself, particularly the clearer handling of state and the separation between parsing and refinement.

the incomplete extension experiment is therefore a limitation of the evaluation. the project can make a strong claim about behavioural preservation, but only a weaker argued claim about maintainability.



\section{Summary}

the evaluation is based primarily on behavioural preservation. the important question is whether the haskell refiner can replace the python one without changing the generated tests or requiring manual adjustment to the workflow. this is a narrower evaluation than originally planned, but it is also the one most closely tied to the main aim of the project.
\chapter{Conclusions \& Future Work}
\label{chap:Conclusions}

This chapter should summarize the work presented in the dissertation and discuss the conclusions that can be drawn from the work and the results presented in chapter~\ref{chap:Evaluation}.


\section{Future Work}

The section may present a list of items that were beyond the scope of the dissertation.

While not as crucial as the porting of the refine-command file,it was originally planned to port spin2test to haskell as well. This would be done as the communication from spin2test to refine command would be vastly simplified, negating the need for the fragile standard input/output type communication that is currently employed. the bridge file and refine command are constantly sending back and forth commands and results. 
in contrast testbuilder launches spin2test once and then lets it read the necessary .spn yaml  and three template files itsel

% \begin{thebibliography}{refs}                   %% Start your bibliography here; you can
\addcontentsline {toc}{chapter}{Bibliography}     %% Force Bibliography to appear in contents
\bibliographystyle{apalike}
\bibliography{refs}                               %% also use the \bibliography command
%\end{thebibliography}                            %% to generate your bibliography.


%\addcontentsline {toc}{chapter}{Appendices}       %% Force Appendices to appear in contents
%\begin{appendix}
%\include{appendix1}
% \include{appendix2}
%\end{appendix}




\end{document}                                    %% END THE DOCUMENT
